\begin{viuabstract}{Abstract}
Cooperative transport of heterogeneous payloads requires a suitable coalition,
physical interaction, and traffic coordination with concurrent missions. This
thesis structures the problem into three subproblems. SP1 decides which robots
serve each payload, how many are needed, which role or contact each member
assumes, and when to wait, start, or switch coalitions. SP2 addresses approach,
docking, stabilisation, transport, and member replacement. SP3 coordinates
routes, reservations, and priorities among free robots, assembling teams, and
moving coalitions. The primary \emph{Cargo} mode models a supported payload as a
planar rigid body.

The proposal combines LP/MILP references, central and neighbour-based sampled
primal--dual dynamics, integer closure with repair, and separate certificates
for resources, contact, and motion. Evaluation uses paired worlds, explicit
seeds, and global oracles as experimental ceilings. The SP1 pilot battery
separates seven questions: a manual case (E0), relaxation (E1), connectivity
(E2), discretisation (E3), integer recovery (E4), scale (E5), and unicycle
arrival (E6).

E0 reproduced the regularised LP in a controlled case, with objective errors on
the order of \(10^{-6}\)--\(10^{-5}\). In E1, Rep-D met the relative stopping
criterion, but only one run was comparable to the LP under the strict primal
tolerance. E2 and E3 remain inconclusive. E4 was feasible in all eight pilot worlds, with a median
gap of \(1.21\%\). In E5, Rep-D repair recovered feasibility at \(N=40\),
although the fractional dynamics did not converge within the iteration budget;
the Greedy baseline left one case infeasible. E6 achieved complete arrival in
two scenarios but did not validate physical transport.

The evidence level is C-pilot: the results verify implementation and controlled
cases, but do not establish global convergence, optimality, or distributed
safety for the complete system. SP2 and SP3 campaigns provide partial evidence;
confirmatory validation, three-dimensional dynamics, wheel--ground traction,
switching networks, hardware, and caging-based transport remain open.

\keywords{Keywords}{multi-robot coordination,
  autonomous mobile robots, coalition formation,
  potential games, cooperative transport}
\end{viuabstract}

\clearpage
